\documentclass[11pt,a4paper]{article}
\usepackage[margin=1.8cm]{geometry}
\usepackage{fontspec}
\setmainfont{Times New Roman}
\setmonofont{Courier New}[Scale=0.85]
\usepackage{polyglossia}
\setdefaultlanguage{vietnamese}
\usepackage{booktabs,tabularx,array,enumitem,fvextra,microtype}
\usepackage[hidelinks]{hyperref}
\setlength{\parindent}{0pt}
\setlength{\parskip}{5pt}
\setlength{\emergencystretch}{2em}
\setlist[itemize]{leftmargin=1.3em,itemsep=2pt,topsep=2pt}
\renewcommand{\arraystretch}{1.15}
\newcolumntype{Y}{>{\raggedright\arraybackslash}X}
\DefineVerbatimEnvironment{cmd}{Verbatim}{fontsize=\small,breaklines=true,breakanywhere=true,frame=single,framesep=4pt}
\newcommand{\heading}[1]{\par\vspace{4pt}{\large\bfseries #1}\par}
\begin{document}
{\LARGE\bfseries EquiCEval: cập nhật mã nguồn}\par
{\small Ngày 15/09/2026 | Đối chiếu với \path{paper_brainstorm/equiceval_v3.tex}}\par
\hrule\vspace{5pt}

\heading{1. Đã bổ sung gì?}
\begin{tabularx}{\textwidth}{@{}p{3.5cm}Y@{}}
\toprule
\textbf{Phần} & \textbf{Cách hoạt động và giới hạn}\\
\midrule
Đối chiếu trong nền chung (M2) & Chỉ dùng các hàng chung đã xác nhận bằng đại số;
loại hàng đang xét và bản sao của chúng. Kiểm tra suy ra hai chiều và kiểm tra
miền nguồn có nghiệm. Tối đa 32 cặp, có thể chưa tìm hết.\\
Đối chiếu nhóm (M2) & Tùy chọn một hàng với hai hàng, tối đa 16 phép thử.
Chỉ xét hàng còn chưa ghép; chọn nhóm theo quy tắc cố định. Không khẳng định
nhóm nhỏ nhất hay phương án ghép nhóm tối ưu.\\
Chỉ số mô tả (M3) & Ghi số điểm thực dùng, số cặp, số phép so sánh, hàng chưa ghép,
hộp và ngưỡng sai số. Không có điểm/cặp thì ghi \texttt{null}, không ghi 0.
Nhóm không được ép thành cặp giả để tính chỉ số.\\
Kiểm tra và mục tiêu & Không gọi miền không bị chặn là lỗi cú pháp; kiểm tra
lại điểm khả thi. Báo riêng hệ số mục tiêu giống nhau hay chỉ cho cùng giá trị.\\
Nhật ký 2.1 & Ghi cặp/nhóm, nền chung, chứng cứ hai chiều, mẫu số và số lần gọi
bộ giải theo bước. Công cụ sinh viên đã đọc định dạng mới.\\
\bottomrule
\end{tabularx}

\heading{2. Ba quy tắc để không kết luận nhầm}
\begin{itemize}
\item \textbf{Đối chiếu thành phần không thay kết luận toàn mô hình.}
Ví dụ, $x+y\leq1$ và $x\leq1$ tương đương khi biết $y=0$, nhưng không tương đương
ở mọi điểm. Phản ví dụ của cặp phải kiểm tra lại toàn bộ nguồn trước khi dùng
để bác bỏ mô hình. Hai miền ngữ cảnh rỗng không được coi là tương đương.
\item \textbf{Phân biệt chứng cứ hữu tỉ với cận số.} \texttt{MatchCoverage}
chỉ tính hàng có chứng cứ hữu tỉ; \texttt{Numerical\allowbreak Match\allowbreak Coverage} nhận thêm
kết luận theo ngưỡng sai số. Mỗi hàng tham chiếu chỉ tính một lần, kể cả trong nhóm.
\item \textbf{M4--M5 chạy trước, giải thích sau.} M2--M3 dùng ngân sách còn lại.
Hết ngân sách vẫn là chưa kết luận, không trở thành ``không có lỗi''.
Điểm mô tả có thể khác nhau ngoài nền chung dù cặp đã được xác nhận trong nền.
\end{itemize}

\heading{3. Chưa được triển khai}
Cross-regret trên \emph{toàn bộ} tập nghiệm tối ưu; tiêu chí chỉ giữ tập nghiệm
tối ưu; đổi đơn vị mục tiêu; chiếu bỏ biến phụ; đối chứng bài gốc đầy đủ;
nhóm giải thích nhỏ nhất. Đầu ra ghi rõ phần chưa hỗ trợ, không sinh số thay thế.

\newpage
{\LARGE\bfseries Kiểm thử và phát hành cấu hình}\par
{\small Dành cho người hướng dẫn; sinh viên chỉ chạy hồ sơ đã được duyệt}\par
\hrule\vspace{5pt}

\heading{4. Kiểm tra đã chạy}
Toàn bộ 89 kiểm thử tự động đã vượt qua, gồm 23 kiểm thử mới so với bản trước.
Nhật ký \path{output/verified_evaluation_20260915_m2m3_final.json} có 108 lượt:
12 ca nhỏ, ba ngân sách và ba chế độ. Ở 0 giây: đều chưa kết luận.
Ở 1 hoặc 5 giây: chế độ có chứng cứ xác nhận 5 ca tương đương, phát hiện 7 ca lỗi;
chế độ chỉ truy vấn xác nhận 5 ca bằng cận số, không gọi là chứng cứ hữu tỉ.
Đây là \textbf{kiểm tra phần mềm}, không phải benchmark LLM thật hay bằng chứng
đủ để khẳng định đóng góp nghiên cứu. Ba chế độ là đối chứng nội bộ, không phải
các phương pháp trong bài báo gốc hoặc các hệ thống khác.

Các kiểm thử mới bao gồm nền $y=0$, thiếu nền, nền rỗng, hàng bị lặp,
nhóm đẳng thức/bất đẳng thức, hết ngân sách, mẫu số, miền không bị chặn và
phát lại chứng cứ hữu tỉ trên hàng gốc. Cảnh báo API PuLP còn tồn tại;
đợt này không nâng cấp thư viện.

\heading{5. Lệnh dành cho thầy, tại thư mục gốc dự án}
\begin{cmd}
.venv/bin/python -m pytest tests -q --disable-warnings
.venv/bin/python -m experiments.run_verified_evaluation --output output/new_check.json --budgets 0 1 5
\end{cmd}
Chọn đường dẫn mới; chương trình không ghi đè kết quả. Mặc định đối chiếu theo
ngữ cảnh bật, nhóm tắt. Khi cần kiểm tra nhóm, dùng lệnh riêng và lưu cấu hình riêng:
\begin{cmd}
.venv/bin/python -m experiments.run_verified_evaluation --output output/new_group_check.json --budgets 5 --group-matching
\end{cmd}
Bộ 12 ca hiện tại không đủ đánh giá lợi ích ghép nhóm; lệnh trên chỉ kiểm tra
đường chạy. Kiểm thử nhóm dùng ca thiết kế riêng. Có thể tắt đối chiếu ngữ cảnh
bằng \texttt{--no-context-matching}; không tự đổi chế độ của sinh viên.

\heading{6. Không dùng hồ sơ cấu hình cũ sau khi sửa mã}
Hồ sơ đã đóng băng trước bản sửa này sẽ bị từ chối do mã băm khác.
Thầy dừng đợt cũ, tạo hồ sơ \emph{mới} trên máy chung và duyệt trước khi giao:
\begin{cmd}
.venv/bin/python -m experiments.controlled_pilot freeze --profile configs/shared_pilot_20260915.json
.venv/bin/python -m experiments.controlled_pilot verify --profile configs/shared_pilot_20260915.json
\end{cmd}
Phiếu chạy của cả Cường, Hoàng, Thành và Sơn phải dùng đúng đường dẫn và
SHA-256 hồ sơ mới. Sinh viên không tự đóng băng, bật ghép nhóm, sửa ngân sách
hay chạy lại khi lỗi. Giữ nguyên kết quả cũ để truy vết; không gộp các phiên bản.
Giống cấu hình vẫn không bảo đảm thời gian giống từng bit hoặc đủ điều kiện làm
bảng hiệu quả chính của paper.

\heading{7. Nơi xem mã}
\path{src/equiceval/matching.py}: nền chung, cặp và nhóm;
\path{src/equiceval/evaluator.py}: đầu ra và chỉ số;
\path{src/equiceval/directed_query.py}: bằng chứng từng chiều;
\path{experiments/run_verified_evaluation.py}: lệnh kiểm tra.
\end{document}
